\begin{figure}[ht]
\centering
\begin{tikzpicture}[>=Stealth, font=\small]
  \draw[very thick, blue!50!black] (0,0) -- (13.8,0);
  \foreach \x/\y in {0.6/1955,2.5/1960s,4.3/1967,6.2/1985,8.0/1986,9.5/1987,11.4/1995,13.0/2001}
    \draw[fill=white] (\x,0) circle (2.2pt) node[below=6pt] {\y};

  \node[above=8pt] at (0.6,0) {谷山提出模性问题};
  \node[above=8pt] at (2.5,0) {志村精确化猜想};
  \node[above=8pt] at (4.3,0) {Weil 提供几何证据};
  \node[above=8pt] at (6.2,0) {Frey 提出反常曲线};
  \node[above=8pt] at (8.0,0) {Ribet 证明水平下降};
  \node[above=8pt] at (9.5,0) {Serre 提出 $\varepsilon$ 猜想};
  \node[above=8pt] at (11.4,0) {Wiles 与 Taylor--Wiles 解决半稳定情形};
  \node[above=8pt] at (13.0,0) {BCDT 完成一般模性};
\end{tikzpicture}
\caption{谷山--志村猜想走向模性定理的时间线：从 1955 年的直觉性问题，到 1995 年的半稳定情形，再到 2001 年的一般情形。}
\label{fig:ch10-history-timeline}
\end{figure}
